\documentstyle[%


righttag%


,11pt%


,amssymb%


,russian%               remove in Eng.


,izvru%                 Set izven% in Eng.


%,izvru-ks             for brief communications


]{amsart}





%





\newcommand{\ctg}{\operatorname{ctg}}


\newcommand{\tg}{\operatorname{tg}}


\newcommand{\cosec}{\operatorname{cosec}}


\newcommand{\gea}{\geqslant}


\newcommand{\lea}{\leqslant}


\newcommand{\vp}{\varphi}


\newcommand{\N}{\Bbb N}


\newcommand{\R}{\Bbb R}


\renewcommand{\L}{\Bbb L}


\newcommand{\intl}{\int\limits}


\newcommand{\suml}{\sum\limits}


\newcommand{\ol}{\overline}


\newcommand{\tts}[1]{}





\theorembodyfont{\rm}


\newtheorem{cornonum}{\hskip\parindent Следствие}


\renewcommand{\thecornonum}{}





%\hoffset=-15mm%                  For Eng.


\setcounter{page}{82}


\god{2003}


\nomer{8~(495)} %                        Remove in Eng.


%\nomer{Vol.\,47, No.\,8}%               Set in Eng.


%\str{\pageref{begin}--\pageref{end}}%  Set in Eng.


\udk{517.968:519.642}





\author{\it И.К.\,Рахимов}


% NB:   ^^ remove in Eng.


\title{Квадратурные методы решения сингулярного интегрального 


уравнения Теодорсена}





\date{}


%\pagestyle{myheadings}%         Set in Eng.


\pagestyle{plain} %              Remove in Eng.


\begin{document}


%\thispagestyle{plain}%          Set in Eng.


\maketitle


\label{begin}











{\it Введение.} Данная статья посвящена квадратурным методам решения 


сингулярного интегрального уравнения (с.\,и.\,у.) Теодорсена 


\begin{equation}


(K\vp)(s)\equiv \vp(s) + \frac{\lambda}{2\pi}\int_{0}^{2\pi}\ln


\rho(\vp(\sigma))\ctg\frac{\sigma-s}{2}d\sigma = y(s),


\end{equation}


где $\lambda$ --- числовой параметр, $\rho(s)$, $y(s)$ ---


известные $2\pi$-периодические функции, а $\vp(s)$ --- искомая функция


и сингулярный интеграл понимается в смысле главного значения по 


Коши--Лебегу (напр., [1]). Это уравнение встречается (напр., [2])


в теории конформных отображений и в ряде прикладных задач. Имеются 


многочисленные результаты по приближенным методам решения уравнения (1)


(напр., [2]--[6], а также [7]). 


\medskip





{\bf 1.} В статье [7] на основе теории монотонных операторов [8] 


доказаны теоремы существования и единственности решения уравнения (1).


Функцию $\vp(s)$ будем искать в вещественном 


пространстве $2\pi$-периодических квад\-ра\-тич\-но-суммируемых


функций $L_2=L_2(0,2\pi)$.


\medskip





{\bf Лемма} ([7], теорема 2). {\it Пусть $\lambda \in \R$ и $\rho(s)$ ---


положительная непрерывно дифференцируемая функция, удовлетворяющая


неравенствам


$$


M\gea \rho'(s)/\rho(s)\gea m >0\ \ \forall s\in\Bbb R.


$$


Тогда уравнение $(1)$ при любой правой части $y(s)\in L_2$ имеет 


единственное решение $\vp^*\in \L_2$.


}


\medskip





{\bf 2.} В приложениях (напр., [9]--[12]) большое значение имеют 


квадратурные методы решения с.\,и.\,у., основанные на аппроксимации 


сингулярных операторов конечномерными операторами, порождаемыми 


квадратурными формулами для интегралов из (1). Рассмотрим три из таких 


методов, основанныx на аппроксимации тригонометрическими 


полиномами и сплайнами.


\medskip





{\bf 2.1.} Введем на $[0,2\pi]$ равноотстоящие узлы


$s_j = \frac{2j\pi}{N}$, $j=\ol{0,N}$, $N\in \N$, и приближенное 


решение уравнения (1) с непрерывными коэффициентами $\ln \rho(s)$


и $y(s)$ будем искать в виде полинома ([9]--[12])


\begin{equation}


\vp_n(s) = \frac{2}{N}\sum_{j=1}^N c_j \Delta_n(s-s_j), \quad n = [N/2],


\end{equation}


где $[t]$ --- целая часть числа $t\gea 0$, а $\Delta_n(s)$ --- 


обычное или же модифицированное ядро Дирихле порядка $n$ ($n+1\in\N$):


$$


\Delta_n(s) = \frac{1}{2}\sin\left(n+\frac{1}{2}\right)


s\cdot \cosec \frac{s}{2},\ \ N = 2n+1;\quad


\Delta_n(s) = \frac{1}{2}\sin ns\cdot \ctg\frac{s}{2},\ \ N = 2n.


$$


Приближенные значения $c_k = \vp_n(s_k)$,


$k=\ol{1,n}$, искомой функции $\vp(s)$ в узлах $s_k$ определяются 


согласно [9]--[12] из системы нелинейных алгебраических уравнений (СНАУ)


\begin{gather}


c_j+\frac{\lambda}{N}\sum_{k=1}^N\ln \rho(c_k)\gamma_{j-k} = y(s_j),


\quad j=\ol{1,N}, \\


%


\gamma_r = \gamma_r^{(N)}=\bigg\{\tg\frac{r\pi}{2N},


\ \; r \text{ четно};\ \; -\ctg\frac{r\pi}{2N},


\ \; r \text{ нечетно}\bigg\},\quad N=2n+1;\notag \\


%


\gamma_r = \gamma_r^{(N)}=\{0,\ \; r \text{ четно };\ \;


-2\ctg\frac{r\pi}{N},\ \; r \text{ нечетно }\},\quad N=2n.\notag


\end{gather}





Далее с.\,и.\,у. (1) будем рассматривать как операторное уравнение 


\begin{gather*}


K\vp \equiv \vp + \lambda I B \vp = y \quad (\vp,y\in L_2), \\


%


I\vp \equiv I(\vp;s) = \frac{1}{2\pi}\int_{0}^{2\pi}\vp(\sigma)


\ctg\frac{\sigma-s}{2}d\sigma,\quad B\vp \equiv \ln \rho(\vp(s)).


\end{gather*}





\begin{thm}


Если в условиях леммы функции $\rho(s)$ и $y(s)$ таковы, что точное 


решение уравнения $(1)$ $\vp(s)\in C_{2\pi}$ и функция $(B\vp)(s)$


удовлетворяет условию Дини--Липшица, то при любых натуральных $n$


система $(3)$ имеет единственное решение и приближенные решения $(2)$


сходятся к точному решению уравнения $(1)$ в $L_2$ со скоростью 


$$


\|\vp - \vp_n\|\lea 2E_{p}(\vp)_C + m^{-1}M|\lambda|\,


\|I(B\vp - \L_nB\vp)\|_C,\ \ n=[N/2],\ \ N\in\N,


$$


где $E_p(\vp)_C$ --- наилучшее равномерное приближение 


функции $\vp(s)\in C_{2\pi}$ тригонометрическими полиномами 


степени $p = [(N-1)/2]$, а


$$


\L_n(f;s) = \frac{2}{N}\suml_{k=1}^N \Delta_n(s-s_k)f(s_k),\quad 


f(s)\in C_{2\pi},


$$


--- оператор Лагранжа.


\end{thm}





\begin{cornonum}


Если в условиях теоремы 1 функции $\rho(s)$


и $y(s)$ таковы, что точное решение уравнения (1)


$\vp(s)\in H_\alpha^{(r)}$, $r\gea 0$, $0<\alpha\lea 1$, то 


приближенные решения (2) сходятся к точному решению уравнения


(1) в $L_2$ и $C_{2\pi}$ со скоростями соответственно


\begin{alignat*}{2}


\|\vp-\vp_n\| &= O(n^{-r-\alpha}),&\quad r+\alpha&>0; \\


%


\|\vp-\vp_n\|_C &= O(n^{-r-\alpha+1/2}),&\quad r+\alpha&>1/2.


\end{alignat*}


\end{cornonum}





{\bf 2.2.} Приближенное решение с.\,и.\,у. (1) ищем в виде сплайна 


нулевого порядка ([10]--[12])


\begin{equation}


\vp_n(s) = \suml_{k=1}^n\alpha_n\chi_k(s),\quad \alpha_k\in\R,\quad


s_k = \frac{2k\pi}{n},\quad 


k=\ol{0,n},\ \ n\in\N,


\end{equation}


где $\chi_k(s)$ --- характеристические функции 


интервалов $(s_{k-1},s_k]$, $k=\ol{1,n}$, а неизвестные 


коэффициенты $\alpha_k=\vp_n(s_k)$ определяются из СНАУ


\begin{gather}


\alpha_j + \lambda\sum_{k=1}^n\ln \rho(\alpha_k)\gamma_{j-k} =


y(\ol{s}_j),\quad j=\ol{1,n},\quad \ol{s}_j = \frac{2j-1}{n}\pi=s_{j-1/2},\\


%


\gamma_{j-k} = \frac{1}{\pi}\ln\left|\frac{\sin s_{k-j+1/2}}


{\sin s_{k-j-1/2}}\right|.


\end{gather}





\begin{thm}


В условиях теоремы {\em 1 СНАУ (5), (6)} однозначно разрешима при 


любых $n\in\N$, а приближенные решения $(4)$ сходятся к точному решению 


со скоростью


$$


\|\vp - \vp_n\|\lea \omega(\vp;2\pi/n)_C +


m^{-1}M|\lambda|\,\|I(B\vp - S_n^0B\vp)\|_C,


$$


где $\omega(\vp;2\pi/n)_C$ --- обычный модуль непрерывности 


функции $\vp(s)\in C_{2\pi}$, а 


$$


S_n^0(f;s) = \sum_{j=1}^n f(\ol{s}_j)\chi_j(s),\quad


f(s)\in C_{2\pi},\quad n\in\N,


$$


причем узлы коллокации $\ol{s}_j$ определены в $(5)$.


\end{thm}





\begin{cornonum}


Пусть в условиях теоремы 2 функции $\rho(s)$ и $y(s)$ таковы, что 


точное решение уравнения (1) $\vp(s)\in H_\alpha^{(r)}$,


$r\gea 0$, $0<\alpha\lea 1$. Тогда метод сплайн-квадратур нулевого 


порядка сходится в среднем и равномерно со скоростями соответственно


\begin{align*}


\|\vp-\vp_n\| &=


\begin{cases}


O(n^{-r-\alpha}),& 0<r+\alpha\lea 1;\\


O(n^{-1}),&r+\alpha>1,


\end{cases}


\\


\|\vp-\vp_n\|_C &=


\begin{cases}


O(n^{-r-\alpha+1/2}),& \frac{1}{2}<r+\alpha\lea 1;\\


O(n^{-1/2}),&r+\alpha>1.


\end{cases}


\end{align*}


\end{cornonum}





{\bf 2.3.} Теперь (напр., [11], [12]) приближенное решение с.\,и.\,у.


(1) будем искать в виде сплайна первой степени


\begin{equation}


\vp_n(s) = \suml_{k=1}^n\alpha_k\phi_k(s),\quad n\in\N,


\end{equation}


где $\{\phi_k(s)\}_1^n$ --- фундаментальные $2\pi$-периодические сплайны 


первой степени для сетки узлов $s_k=\frac{2k\pi}{n}$,


$k=\ov{0,n}:\phi_k(s) = \{0$ при


$s\lea s_{k-1}$ и $s\gea s_{k+1}$;


$(s-s_{k-1})/\delta$ при $s_{k-1}\lea s\lea s_k$;


$(s_{k+1}-s)/\delta$ при $s_k\lea s\lea s_{k+1}\}$, $\delta =2\pi/n$. 


Неизвестные коэффициенты $\alpha_k$ определим ([11], с. 152) из СНАУ


\begin{gather}


\alpha_j + \lambda\sum_{k=1}^n \ln \rho(\alpha_k)\gamma_{j-k} = y(s_{j}),


\quad j=\ol{1,n}, \\


%


\gamma_{j-k} = \frac{1}{2\pi}\int_{0}^\delta \left[


\ctg\frac{t-(j-k)\delta}{2}-\ctg\frac{t-(k-j)\delta}{2}


\right]\left(1-\frac{t}{\delta}\right)dt,\quad \delta = \frac{2\pi}{n}.


\end{gather}





\begin{thm}


В условиях теоремы {\em 1 СНАУ (8), (9)} однозначно разрешима 


при любых $n\in\N$, а приближенные решения $(7)$ сходятся к точному 


решению со скоростью 


$$


\|\vp - \vp_n\|\lea \omega(\vp;2\pi/n)_C +


\sqrt{3}m^{-1}M|\lambda|\,\|I(B\vp - S_n^1B\vp)\|_C,


$$


где 


$$


S_n^1(f;s) = \sum_{j=1}^n f(s_j)\phi_j(s),


\quad f(s)\in C_{2\pi},\quad n\in\N.


$$


\end{thm}





\begin{cornonum}


Пусть в условиях теоремы 3 функции $\rho(s)$ и $y(s)$ таковы, что 


точное решение уравнения (1) $\vp(s)\in H^{(r)}_\alpha$,


$r\gea 0$, $0<\alpha\lea 1$. Тогда метод сплайн-квадратур первого 


порядка сходится в среднем и равномерно со скоростями соответственно


\begin{align*}


\|\vp - \vp_n\| &= 


\begin{cases}


O(n^{-r-\alpha}),&0<r+\alpha\lea 2; \\


O(n^{-2}),&r+\alpha>2,


\end{cases}


\\


\|\vp - \vp_n\|_C &=


\begin{cases}


O(n^{-r-\alpha+1/2}),&\frac{1}{2}<r+\alpha\lea 2; \\


O(n^{-3/2}),&r+\alpha>2.


\end{cases}


\end{align*}


\end{cornonum}





При обосновании указанных квадратурных методов были использованы 


результаты работ [6]--[14].





\begin{note}


Теоремы 1--3 и результаты ([9], \S\,7, гл.\,2) позволяют построить 


сходящиеся в $L_2(0,2\pi)$ схемы квадратурно-итерационных методов 


решения с.\,и.\,у. (1).


\end{note}





\begin{note}


В условиях теоремы 1 из [7] справедливы результаты, аналогичные 


полученным выше.


\end{note}





\begin{thebibliography}{99}





\bibitem{1}


Michlin S.G., Pr\"o\ss{}dorf S. {\em Singul\"{a}re Integraloperatoren}.


-- Berlin: Akademie-Verlag, 1980. -- 514 S.


\bibitem{2}


Gaier D. {\em Numerical methods in conformal mapping} // Comput. Aspects 


Complex Anal. Proc. NATO Adv. Stady Inst., Braunlage, Harz, 26 July -- 6


Aug. 1982. -- 1983. -- P.\,51--78.


\bibitem{3}


Габдулхаев Б.Г. {\em Конечномерные аппроксимации сингулярных интегралов


и прямые методы решения особых интегральных и интегро-дифференциальных


уравнений} // Итоги науки и техн. ВИНИТИ. Матем. анализ. -- 1980. --


Т.\,18. -- С.\,251--307.


\bibitem{4}


Мусаев Б.И. {\em Конструктивные методы в теории сингулярных интегральных


уравнений}: Дис. $\dotsc$ докт. физ.-матем. наук. -- Тбилиси, 1989. -- 339 с.


\bibitem{5}


Лифанов И.К., Тыртышников Е.Е. {\em Теплицeвы матрицы и сингулярные 


интегральные уравнения} // Вычисл. процессы и системы. -- М.: Наука, 


1990. -- Вып.\,7. -- С.\,94--278.


\bibitem{6}


Рахимов И.К. {\em Прямые методы решения нелинейных сингулярных 


интегральных уравнений с монотонными операторами}: Дис. $\dotsc$


канд. физ.-матем. наук. -- Казань, 1998. -- 150 с.


\bibitem{7}


Рахимов И.К. {\em Проекционные методы решения сингулярного интегрального 


уравнения Теодорсена} // Изв. вузов. Математика. -- 2002.


-- \No\,9. -- С.\,67--70.


\bibitem{8}


Гаевский Х., Грег\"eр К., Захариас К. {\em Нелинейные операторные 


уравнения и операторные дифференциальные уравнения}. -- М.: Мир,


1978. -- 336 с.


\bibitem{9}


Габдулхаев Б.Г. {\em Оптимальные аппроксимации решений линейных задач}.


-- Казань: Изд-во Казанск. ун-та, 1980. -- 232 с.


\bibitem{10}


Габдулхаев Б.Г. {\em Прямые методы решения сингулярных интегральных 


уравнений первого рода. Численный анализ}. -- Казань: Изд-во Казанск.


ун-та, 1994. -- 288 с.


\bibitem{11}


Габдулхаев Б.Г. {\em Численный анализ сингулярных интегральных уравнений. 


Избранные главы}. -- Казань: Изд-во Казанск. ун-та, 1995. -- 230 с.


\bibitem{12}


Габдулхаев Б.Г. {\em Методы решения сингулярных интегральных уравнений 


с положительными операторами} // Дифференц. уравнения.


-- 1997. -- Т.\,33. -- \No\,3. -- С.\,400--409.


\bibitem{13}


Корнейчук Н.П. {\em Точные константы в теории приближения}.


-- М.: Наука, 1987. -- 423 с.


\bibitem{4}


Корнейчук Н.П. {\em Сплайны в теории приближения}.


-- М.: Наука, 1984. -- 352 с.





\end{thebibliography}





\vskip 2cm





\postup{\it Казанская государственная}{\qquad \it Поступили}


\postup{\it сельскохозяйственная академия}{\it полный текст $11.02.2002$}


\postup{}{\it краткое сообщение $13.03.2003$}





\label{end}


\end{document}


